\section{语料与口径}

\subsection{语料与抽取产物}

本文语料是通俗明史叙事《明朝那些事儿》全七部（2006--2009 年间连载出版）\citep{dangnian2009}，以及在其之上抽取的结构化知识图谱。图谱不是原书本身，而是原书的一次投影；本文全部论据都建立在这一投影之上，第 6 节会反复回到这一点。

规模如表 \ref{tab:corpus} 所示。人物 1231、事件 1106、关系 2324、地点 563；此外还有作者标注的在位表与人物年谱两项辅助表，本文不使用。

\begin{table}[!t]
\centering
\caption{语料与结构化抽取产物规模}
\label{tab:corpus}
\small
\begin{tabular}{lr}
\toprule
项目 & 数量 \\
\midrule
原书全书章节 & 156 \\
本篇锚点覆盖章节 & 41 \\
人物 & 1231 \\
事件 & 1106 \\
关系边 & 2324 \\
地点 & 563 \\
\bottomrule
\end{tabular}
\end{table}



\subsection{主题词表的构造}

本文关心的是「家国同构——继承、名分与亲属称谓」这组词，词表由 12 个中文词构成：族、孝、嗣、继承、家族、名分、世袭、祭祀、嫡子、荫、宗族、家法。命中一律按字符逐字在原书文本中进行，不做同义词扩展、不做分词、不把现代术语换算进去。词表共 14 个候选词，逐字命中原书后得到 12 个不同词进表（另有「嫡长」「宗法」两词
零命中，按规则剔除，不计入任何合计）；命中不做同义词扩展、不做分词、不把现代术语换算进去。
注意短字（族、孝、嗣、荫）一字命中会带上同字根的无关词，这一条的后果写在第 6 节 L2。

\begin{table}[!t]
\centering
\caption{主题词逐字命中（降序）}
\label{tab:vocab}
\small
\begin{tabular}{lr}
\toprule
词 & 命中次数 \\
\midrule
族 & 234 \\
孝 & 167 \\
嗣 & 114 \\
继承 & 82 \\
家族 & 37 \\
名分 & 34 \\
世袭 & 23 \\
祭祀 & 18 \\
嫡子 & 17 \\
荫 & 7 \\
宗族 & 1 \\
家法 & 1 \\
\bottomrule
\end{tabular}
\end{table}



\subsection{锚点}

为可核验起见另列 4 个锚点，全部取自人物表且原书确有此人。度数与出场章节数如表 \ref{tab:anchors}。

\begin{table}[!t]
\centering
\caption{锚点人物的度数与出场章数（按出场章数降序，取前 8）}
\label{tab:anchors}
\small
\begin{tabular}{lrr}
\toprule
人物 & 度数 & 出场章数 \\
\midrule
朱元璋 & 87 & 41 \\
朱棣 & 113 & 25 \\
朱由检 & 67 & 19 \\
方孝孺 & 15 & 6 \\
\bottomrule
\end{tabular}
\end{table}

